% Seminar 3 - Pagina 1, manuscris datat 24-oct.-1994.
% Fragment pentru % Seminar 3 - Pagina 1, manuscris datat 24-oct.-1994.
% Fragment pentru % Seminar 3 - Pagina 1, manuscris datat 24-oct.-1994.
% Fragment pentru % Seminar 3 - Pagina 1, manuscris datat 24-oct.-1994.
% Fragment pentru \input{seminars/pages/seminar3-pagina1-ro}; fără preambul.
% NOTE EDITORIALE:
% - Problema 1 reia deranjamentele de la Seminarul 2.
%   La început notația seamănă cu delta(n), iar în ultima formulă
%   apare D_n. Sunt redate ambele notații în pozițiile respective.
% - Demonstrația prin principiul includerii și excluderii este redată
%   în ordinea manuscrisului, inclusiv termenul pentru k=n.
% - În problema 2(a), o parte din expresia inițială Hom este tăiată;
%   este folosită expresia finală lizibilă Hom(X x Y, Z).
% - Aplicația inversă a lui alpha este numită beta în demonstrație;
%   beta a fost folosit și pentru cealaltă corespondență pe pagina
%   precedentă a Seminarului 2. Denumirea locală este păstrată.
% - Au fost omise doar calcule și simboluri barate, greu lizibile.
% - Manuscrisul afirmă corespondența de la (b), dar demonstrează
%   numai punctul (a) pe această pagină; nu am adăugat demonstrația (b).

% \phantomsection
\section{Seminar 3}
% \addcontentsline{toc}{subsection}{Seminar 3}
\begin{flushright}
\textit{24-oct.-1994}
\end{flushright}

\noindent\fbox{1}\quad O permutare $\sigma\in S_n$ se numeşte
\emph{deranjament} dacă
\[
  \sigma(i)\neq i,\qquad \forall i=1,\ldots,n.
\]
Fie $n\in\mathbb N^{*}$; notăm
\[
 \delta(n)=\bigl|\{\sigma\in S_n\mid\sigma\text{ este deranjament}\}\bigr|.
\]
Să se calculeze $\delta(n)$.

\medskip
\noindent Fie
\[
 A_i=\{\sigma\in S_n\mid\sigma(i)=i\},
 \qquad |A_i|=(n-1)!.
\]
Atunci
\[
 \delta(n)=n!-\bigl|\{\sigma\in S_n\mid
                   \sigma\text{ nu este deranjament}\}\bigr|.
\]
Notăm cu $B$ mulţimea dintre acolade. Deci
\[
 B=A_1\cup\cdots\cup A_n.
\]
Din principiul includerii şi excluderii,
\[
\begin{aligned}
 |B|
 &=|A_1\cup\cdots\cup A_n|\\
 &=\sum_{i=1}^{n}|A_i|
    -\sum_{1\le i<j\le n}|A_i\cap A_j|+\cdots
      +(-1)^{n-1}|A_1\cap\cdots\cap A_n|\\[3pt]
 &=\binom n1(n-1)!-\binom n2(n-2)!+\cdots
      +(-1)^{n-1}\binom nn\,0!\\[3pt]
 &=\sum_{k=1}^{n}(-1)^{k-1}\binom nk(n-k)!.
\end{aligned}
\]
Prin urmare, în notaţia de la sfârşitul manuscrisului,
\[
 \boxed{D_n=n!-|B|=\sum_{k=0}^{n}(-1)^k\binom nk(n-k)!}.
\]

\bigskip
\noindent\fbox{2}\quad Fie $X,Y,Z$ mulţimi.

\medskip
\noindent\textbf{(a)} Avem corespondenţa:
\[
 \operatorname{Hom}(X\times Y,Z)
 \;\overset{\alpha}{\underset{\beta}{\rightleftarrows}}\;
 \operatorname{Hom}\bigl(X,\operatorname{Hom}(Y,Z)\bigr),
\]
unde
\[
 \alpha(f)(x)(y)=f(x,y).
\]

\noindent\textbf{(b)} De asemenea,
\[
 \operatorname{Hom}(X,Y\times Z)
 \cong\operatorname{Hom}(X,Y)\times\operatorname{Hom}(X,Z).
\]

\medskip
\noindent\underline{Demonstraţia punctului (a).}
Definim aplicaţia inversă prin
\[
 \beta(g)(a,b)=g(a)(b),\qquad
 g:X\longrightarrow\operatorname{Hom}(Y,Z).
\]
Pentru orice astfel de $g$,
\[
 (\alpha\circ\beta)(g)=g.
\]
Într-adevăr, pentru toţi $x\in X$ şi $y\in Y$,
\[
 \begin{aligned}
 (\alpha\circ\beta)(g)(x)(y)
 &=\alpha\bigl(\beta(g)\bigr)(x)(y)\\
 &=\beta(g)(x,y)\\
 &=g(x)(y).
 \end{aligned}
\]
Pe de altă parte, pentru $f:X\times Y\longrightarrow Z$,
\[
 (\beta\circ\alpha)(f)=\operatorname{id}_{\operatorname{Hom}(X\times Y,Z)}(f)=f.
\]
Într-adevăr, pentru orice $x\in X$ şi $y\in Y$,
\[
 \begin{aligned}
 \beta\bigl(\alpha(f)\bigr)(x,y)
 &=\alpha(f)(x)(y)\\
 &=f(x,y).
 \end{aligned}
\]
; fără preambul.
% NOTE EDITORIALE:
% - Problema 1 reia deranjamentele de la Seminarul 2.
%   La început notația seamănă cu delta(n), iar în ultima formulă
%   apare D_n. Sunt redate ambele notații în pozițiile respective.
% - Demonstrația prin principiul includerii și excluderii este redată
%   în ordinea manuscrisului, inclusiv termenul pentru k=n.
% - În problema 2(a), o parte din expresia inițială Hom este tăiată;
%   este folosită expresia finală lizibilă Hom(X x Y, Z).
% - Aplicația inversă a lui alpha este numită beta în demonstrație;
%   beta a fost folosit și pentru cealaltă corespondență pe pagina
%   precedentă a Seminarului 2. Denumirea locală este păstrată.
% - Au fost omise doar calcule și simboluri barate, greu lizibile.
% - Manuscrisul afirmă corespondența de la (b), dar demonstrează
%   numai punctul (a) pe această pagină; nu am adăugat demonstrația (b).

% \phantomsection
\section{Seminar 3}
% \addcontentsline{toc}{subsection}{Seminar 3}
\begin{flushright}
\textit{24-oct.-1994}
\end{flushright}

\noindent\fbox{1}\quad O permutare $\sigma\in S_n$ se numeşte
\emph{deranjament} dacă
\[
  \sigma(i)\neq i,\qquad \forall i=1,\ldots,n.
\]
Fie $n\in\mathbb N^{*}$; notăm
\[
 \delta(n)=\bigl|\{\sigma\in S_n\mid\sigma\text{ este deranjament}\}\bigr|.
\]
Să se calculeze $\delta(n)$.

\medskip
\noindent Fie
\[
 A_i=\{\sigma\in S_n\mid\sigma(i)=i\},
 \qquad |A_i|=(n-1)!.
\]
Atunci
\[
 \delta(n)=n!-\bigl|\{\sigma\in S_n\mid
                   \sigma\text{ nu este deranjament}\}\bigr|.
\]
Notăm cu $B$ mulţimea dintre acolade. Deci
\[
 B=A_1\cup\cdots\cup A_n.
\]
Din principiul includerii şi excluderii,
\[
\begin{aligned}
 |B|
 &=|A_1\cup\cdots\cup A_n|\\
 &=\sum_{i=1}^{n}|A_i|
    -\sum_{1\le i<j\le n}|A_i\cap A_j|+\cdots
      +(-1)^{n-1}|A_1\cap\cdots\cap A_n|\\[3pt]
 &=\binom n1(n-1)!-\binom n2(n-2)!+\cdots
      +(-1)^{n-1}\binom nn\,0!\\[3pt]
 &=\sum_{k=1}^{n}(-1)^{k-1}\binom nk(n-k)!.
\end{aligned}
\]
Prin urmare, în notaţia de la sfârşitul manuscrisului,
\[
 \boxed{D_n=n!-|B|=\sum_{k=0}^{n}(-1)^k\binom nk(n-k)!}.
\]

\bigskip
\noindent\fbox{2}\quad Fie $X,Y,Z$ mulţimi.

\medskip
\noindent\textbf{(a)} Avem corespondenţa:
\[
 \operatorname{Hom}(X\times Y,Z)
 \;\overset{\alpha}{\underset{\beta}{\rightleftarrows}}\;
 \operatorname{Hom}\bigl(X,\operatorname{Hom}(Y,Z)\bigr),
\]
unde
\[
 \alpha(f)(x)(y)=f(x,y).
\]

\noindent\textbf{(b)} De asemenea,
\[
 \operatorname{Hom}(X,Y\times Z)
 \cong\operatorname{Hom}(X,Y)\times\operatorname{Hom}(X,Z).
\]

\medskip
\noindent\underline{Demonstraţia punctului (a).}
Definim aplicaţia inversă prin
\[
 \beta(g)(a,b)=g(a)(b),\qquad
 g:X\longrightarrow\operatorname{Hom}(Y,Z).
\]
Pentru orice astfel de $g$,
\[
 (\alpha\circ\beta)(g)=g.
\]
Într-adevăr, pentru toţi $x\in X$ şi $y\in Y$,
\[
 \begin{aligned}
 (\alpha\circ\beta)(g)(x)(y)
 &=\alpha\bigl(\beta(g)\bigr)(x)(y)\\
 &=\beta(g)(x,y)\\
 &=g(x)(y).
 \end{aligned}
\]
Pe de altă parte, pentru $f:X\times Y\longrightarrow Z$,
\[
 (\beta\circ\alpha)(f)=\operatorname{id}_{\operatorname{Hom}(X\times Y,Z)}(f)=f.
\]
Într-adevăr, pentru orice $x\in X$ şi $y\in Y$,
\[
 \begin{aligned}
 \beta\bigl(\alpha(f)\bigr)(x,y)
 &=\alpha(f)(x)(y)\\
 &=f(x,y).
 \end{aligned}
\]
; fără preambul.
% NOTE EDITORIALE:
% - Problema 1 reia deranjamentele de la Seminarul 2.
%   La început notația seamănă cu delta(n), iar în ultima formulă
%   apare D_n. Sunt redate ambele notații în pozițiile respective.
% - Demonstrația prin principiul includerii și excluderii este redată
%   în ordinea manuscrisului, inclusiv termenul pentru k=n.
% - În problema 2(a), o parte din expresia inițială Hom este tăiată;
%   este folosită expresia finală lizibilă Hom(X x Y, Z).
% - Aplicația inversă a lui alpha este numită beta în demonstrație;
%   beta a fost folosit și pentru cealaltă corespondență pe pagina
%   precedentă a Seminarului 2. Denumirea locală este păstrată.
% - Au fost omise doar calcule și simboluri barate, greu lizibile.
% - Manuscrisul afirmă corespondența de la (b), dar demonstrează
%   numai punctul (a) pe această pagină; nu am adăugat demonstrația (b).

% \phantomsection
\section{Seminar 3}
% \addcontentsline{toc}{subsection}{Seminar 3}
\begin{flushright}
\textit{24-oct.-1994}
\end{flushright}

\noindent\fbox{1}\quad O permutare $\sigma\in S_n$ se numeşte
\emph{deranjament} dacă
\[
  \sigma(i)\neq i,\qquad \forall i=1,\ldots,n.
\]
Fie $n\in\mathbb N^{*}$; notăm
\[
 \delta(n)=\bigl|\{\sigma\in S_n\mid\sigma\text{ este deranjament}\}\bigr|.
\]
Să se calculeze $\delta(n)$.

\medskip
\noindent Fie
\[
 A_i=\{\sigma\in S_n\mid\sigma(i)=i\},
 \qquad |A_i|=(n-1)!.
\]
Atunci
\[
 \delta(n)=n!-\bigl|\{\sigma\in S_n\mid
                   \sigma\text{ nu este deranjament}\}\bigr|.
\]
Notăm cu $B$ mulţimea dintre acolade. Deci
\[
 B=A_1\cup\cdots\cup A_n.
\]
Din principiul includerii şi excluderii,
\[
\begin{aligned}
 |B|
 &=|A_1\cup\cdots\cup A_n|\\
 &=\sum_{i=1}^{n}|A_i|
    -\sum_{1\le i<j\le n}|A_i\cap A_j|+\cdots
      +(-1)^{n-1}|A_1\cap\cdots\cap A_n|\\[3pt]
 &=\binom n1(n-1)!-\binom n2(n-2)!+\cdots
      +(-1)^{n-1}\binom nn\,0!\\[3pt]
 &=\sum_{k=1}^{n}(-1)^{k-1}\binom nk(n-k)!.
\end{aligned}
\]
Prin urmare, în notaţia de la sfârşitul manuscrisului,
\[
 \boxed{D_n=n!-|B|=\sum_{k=0}^{n}(-1)^k\binom nk(n-k)!}.
\]

\bigskip
\noindent\fbox{2}\quad Fie $X,Y,Z$ mulţimi.

\medskip
\noindent\textbf{(a)} Avem corespondenţa:
\[
 \operatorname{Hom}(X\times Y,Z)
 \;\overset{\alpha}{\underset{\beta}{\rightleftarrows}}\;
 \operatorname{Hom}\bigl(X,\operatorname{Hom}(Y,Z)\bigr),
\]
unde
\[
 \alpha(f)(x)(y)=f(x,y).
\]

\noindent\textbf{(b)} De asemenea,
\[
 \operatorname{Hom}(X,Y\times Z)
 \cong\operatorname{Hom}(X,Y)\times\operatorname{Hom}(X,Z).
\]

\medskip
\noindent\underline{Demonstraţia punctului (a).}
Definim aplicaţia inversă prin
\[
 \beta(g)(a,b)=g(a)(b),\qquad
 g:X\longrightarrow\operatorname{Hom}(Y,Z).
\]
Pentru orice astfel de $g$,
\[
 (\alpha\circ\beta)(g)=g.
\]
Într-adevăr, pentru toţi $x\in X$ şi $y\in Y$,
\[
 \begin{aligned}
 (\alpha\circ\beta)(g)(x)(y)
 &=\alpha\bigl(\beta(g)\bigr)(x)(y)\\
 &=\beta(g)(x,y)\\
 &=g(x)(y).
 \end{aligned}
\]
Pe de altă parte, pentru $f:X\times Y\longrightarrow Z$,
\[
 (\beta\circ\alpha)(f)=\operatorname{id}_{\operatorname{Hom}(X\times Y,Z)}(f)=f.
\]
Într-adevăr, pentru orice $x\in X$ şi $y\in Y$,
\[
 \begin{aligned}
 \beta\bigl(\alpha(f)\bigr)(x,y)
 &=\alpha(f)(x)(y)\\
 &=f(x,y).
 \end{aligned}
\]
; fără preambul.
% NOTE EDITORIALE:
% - Problema 1 reia deranjamentele de la Seminarul 2.
%   La început notația seamănă cu delta(n), iar în ultima formulă
%   apare D_n. Sunt redate ambele notații în pozițiile respective.
% - Demonstrația prin principiul includerii și excluderii este redată
%   în ordinea manuscrisului, inclusiv termenul pentru k=n.
% - În problema 2(a), o parte din expresia inițială Hom este tăiată;
%   este folosită expresia finală lizibilă Hom(X x Y, Z).
% - Aplicația inversă a lui alpha este numită beta în demonstrație;
%   beta a fost folosit și pentru cealaltă corespondență pe pagina
%   precedentă a Seminarului 2. Denumirea locală este păstrată.
% - Au fost omise doar calcule și simboluri barate, greu lizibile.
% - Manuscrisul afirmă corespondența de la (b), dar demonstrează
%   numai punctul (a) pe această pagină; nu am adăugat demonstrația (b).

% \phantomsection
\section{Seminar 3}
% \addcontentsline{toc}{subsection}{Seminar 3}
\begin{flushright}
\textit{24-oct.-1994}
\end{flushright}

\noindent\fbox{1}\quad O permutare $\sigma\in S_n$ se numeşte
\emph{deranjament} dacă
\[
  \sigma(i)\neq i,\qquad \forall i=1,\ldots,n.
\]
Fie $n\in\mathbb N^{*}$; notăm
\[
 \delta(n)=\bigl|\{\sigma\in S_n\mid\sigma\text{ este deranjament}\}\bigr|.
\]
Să se calculeze $\delta(n)$.

\medskip
\noindent Fie
\[
 A_i=\{\sigma\in S_n\mid\sigma(i)=i\},
 \qquad |A_i|=(n-1)!.
\]
Atunci
\[
 \delta(n)=n!-\bigl|\{\sigma\in S_n\mid
                   \sigma\text{ nu este deranjament}\}\bigr|.
\]
Notăm cu $B$ mulţimea dintre acolade. Deci
\[
 B=A_1\cup\cdots\cup A_n.
\]
Din principiul includerii şi excluderii,
\[
\begin{aligned}
 |B|
 &=|A_1\cup\cdots\cup A_n|\\
 &=\sum_{i=1}^{n}|A_i|
    -\sum_{1\le i<j\le n}|A_i\cap A_j|+\cdots
      +(-1)^{n-1}|A_1\cap\cdots\cap A_n|\\[3pt]
 &=\binom n1(n-1)!-\binom n2(n-2)!+\cdots
      +(-1)^{n-1}\binom nn\,0!\\[3pt]
 &=\sum_{k=1}^{n}(-1)^{k-1}\binom nk(n-k)!.
\end{aligned}
\]
Prin urmare, în notaţia de la sfârşitul manuscrisului,
\[
 \boxed{D_n=n!-|B|=\sum_{k=0}^{n}(-1)^k\binom nk(n-k)!}.
\]

\bigskip
\noindent\fbox{2}\quad Fie $X,Y,Z$ mulţimi.

\medskip
\noindent\textbf{(a)} Avem corespondenţa:
\[
 \operatorname{Hom}(X\times Y,Z)
 \;\overset{\alpha}{\underset{\beta}{\rightleftarrows}}\;
 \operatorname{Hom}\bigl(X,\operatorname{Hom}(Y,Z)\bigr),
\]
unde
\[
 \alpha(f)(x)(y)=f(x,y).
\]

\noindent\textbf{(b)} De asemenea,
\[
 \operatorname{Hom}(X,Y\times Z)
 \cong\operatorname{Hom}(X,Y)\times\operatorname{Hom}(X,Z).
\]

\medskip
\noindent\underline{Demonstraţia punctului (a).}
Definim aplicaţia inversă prin
\[
 \beta(g)(a,b)=g(a)(b),\qquad
 g:X\longrightarrow\operatorname{Hom}(Y,Z).
\]
Pentru orice astfel de $g$,
\[
 (\alpha\circ\beta)(g)=g.
\]
Într-adevăr, pentru toţi $x\in X$ şi $y\in Y$,
\[
 \begin{aligned}
 (\alpha\circ\beta)(g)(x)(y)
 &=\alpha\bigl(\beta(g)\bigr)(x)(y)\\
 &=\beta(g)(x,y)\\
 &=g(x)(y).
 \end{aligned}
\]
Pe de altă parte, pentru $f:X\times Y\longrightarrow Z$,
\[
 (\beta\circ\alpha)(f)=\operatorname{id}_{\operatorname{Hom}(X\times Y,Z)}(f)=f.
\]
Într-adevăr, pentru orice $x\in X$ şi $y\in Y$,
\[
 \begin{aligned}
 \beta\bigl(\alpha(f)\bigr)(x,y)
 &=\alpha(f)(x)(y)\\
 &=f(x,y).
 \end{aligned}
\]
